\documentclass[11pt,a4paper]{article}

\usepackage{../shared/cathedral-preamble}
\usepackage{setspace}
\onehalfspacing

\title{%
  The Cathedral: A Guide for the Curious\\[4pt]
  \large What a Computer Proof About Prime Numbers\\
  \large Means for the Rest of Us%
}
\author{Jason Robert Gochanour}
\date{April 28, 2026}

\begin{document}
\maketitle

\begin{abstract}
This paper explains the Cathedral project---a computer-verified
proof about the Riemann Hypothesis---for readers without a
mathematics background. No equations are required. We explain what
the Riemann Hypothesis is, why it matters, what the Cathedral
actually proved, and why the answer is ``we reduced a 167-year-old
mystery to two well-understood facts, and a computer checked
our work.''
\end{abstract}

% ================================================================
\section{The Oldest Open Problem}
% ================================================================

There is a question about prime numbers---2, 3, 5, 7, 11, 13,
\ldots---that has been open since 1859. A German mathematician
named Bernhard Riemann conjectured a pattern in how primes are
distributed among the whole numbers. His conjecture, called the
\emph{Riemann Hypothesis}, has never been proved or disproved.

It is considered the most important unsolved problem in mathematics.
A proof would have implications for cryptography, physics, and our
understanding of the fundamental structure of numbers.

% ================================================================
\section{What We Did}
% ================================================================

We did \textbf{not} prove the Riemann Hypothesis.

What we did is this: we built a formal proof---checked by a
computer, line by line---that the Riemann Hypothesis is
\emph{exactly equivalent} to a different mathematical statement.
This other statement says that a certain distance measurement
(called the Nyman--Beurling distance) gets closer and closer to
zero.

Think of it like this: imagine you want to know whether a mountain
has gold inside it. You can't dig through the mountain. But you
\emph{can} prove, with mathematical certainty, that the mountain
has gold \emph{if and only if} a certain needle on a gauge points
to zero. You haven't found the gold, but you've built a perfect
measuring instrument.

% ================================================================
\section{What Does ``Computer-Verified'' Mean?}
% ================================================================

When a human writes a mathematical proof, other humans check it.
Humans make mistakes. In 2012, a proof of a major conjecture was
published, celebrated, and later found to contain an error.

When a computer verifies a proof, it checks every logical step
against the rules of mathematics, mechanically, without skipping
anything. It cannot be fooled, bribed, or tired. If the computer
says the proof is correct, it is correct---up to the correctness
of the computer itself.

We used a program called \textbf{Lean~4}, developed at Microsoft
Research, which has been used to verify proofs in algebra, geometry,
and analysis. Our proof is 42{,}491 lines long and compiles without
errors.

% ================================================================
\section{What Remains}
% ================================================================

Our proof reduces the Riemann Hypothesis to \textbf{two} remaining
facts. Both are well-known results that mathematicians have
rigorously proven. They are:

\begin{enumerate}
  \item A bound on how the Riemann zeta function behaves on a
    certain line (proved by Hardy and Littlewood in 1918).
  \item A bound on how large the zeta function can be in a certain
    region (proved by Hadamard in 1893).
\end{enumerate}

Both of these are accepted mathematics. They are ``axioms'' in our
system only because the computer's math library doesn't yet contain
the infrastructure to verify them mechanically. When it does, our
proof will be complete.

% ================================================================
\section{The Three Builders}
% ================================================================

The Cathedral was built by three minds working together:

\begin{enumerate}
  \item A human mathematician (Jason Robert Gochanour), who provided
    the vision, selected the approach, and checked every line.
  \item Google's Gemini AI, which suggested mathematical strategies
    and identified key architectural decisions.
  \item Anthropic's Claude AI, which wrote the computer-verified
    proofs and built the numerical experiments.
\end{enumerate}

This is, to our knowledge, the first time two AI systems and a
human have collaborated to produce a formally verified mathematical
result of this scale.

% ================================================================
\section{Why Does This Matter?}
% ================================================================

\begin{enumerate}
  \item \textbf{For mathematics}: It shows that computer verification
    can tackle the hardest open problems, not just textbook exercises.
  \item \textbf{For AI}: It demonstrates a model where AI augments
    human capability rather than replacing it.
  \item \textbf{For everyone}: It means the logical structure of one
    of humanity's deepest mathematical mysteries is now mapped,
    verified, and publicly available for anyone to inspect.
\end{enumerate}

% ================================================================
\section{For Journalists}
% ================================================================

\subsection{One-Sentence Summary}

A human-AI team built a computer-verified proof that reduces the
167-year-old Riemann Hypothesis to two well-understood mathematical
facts.

\subsection{Key Facts}

\begin{itemize}
  \item \textbf{What}: Formal proof that RH is equivalent to a
    distance measurement going to zero.
  \item \textbf{Who}: Jason Robert Gochanour with Google's Gemini
    and Anthropic's Claude.
  \item \textbf{How long}: 32 days, March--April 2026.
  \item \textbf{Size}: 42{,}491 lines of verified proof code,
    38 computational experiments.
  \item \textbf{What it is NOT}: A proof of the Riemann Hypothesis.
  \item \textbf{What it IS}: A machine-verified reduction to two
    standard analytic bounds.
\end{itemize}

\subsection{Suggested Headline}

\emph{``Computer Proof Maps the Logical Structure of Math's Greatest
Unsolved Problem''}

\subsection{Common Misconceptions}

\begin{enumerate}
  \item \textbf{``They proved RH!''} --- No. We proved an equivalence.
  \item \textbf{``AI solved a math problem.''} --- AI assisted; the
    human directed and verified.
  \item \textbf{``This is just a computer program.''} --- It is a
    proof, checked by a theorem-proving compiler.
\end{enumerate}

% ================================================================
\section{How to Verify}
% ================================================================

The entire proof is open source. To check it yourself:

\begin{enumerate}
  \item Install Lean~4 (\url{https://leanprover.github.io}).
  \item Clone the repository.
  \item Run \texttt{cd proofs \&\& lake build}.
  \item If it compiles: the proof is correct.
\end{enumerate}

No trust in the authors is required. The compiler is the arbiter.

\begin{thebibliography}{99}
\bibitem{cathedral}
J.R.~Gochanour, \emph{The Cathedral}, 2026.
\end{thebibliography}

\end{document}
